\chapter{Introduction}
\label{ch:intro}

Modern software systems are becoming increasingly complex, often consisting of
multiple components written in different programming languages. This diversity
is driven by the need to leverage the strengths of various technologies,
frameworks, and ecosystems. As a consequence, developers frequently face the
challenge of translating or migrating code between programming languages to
ensure maintainability, scalability, and long-term sustainability of software
systems.

In recent years, this problem has increasingly been addressed using deep
learning, particularly Large Language Models (LLMs). These models are capable
of understanding programming languages as structured text and can automatically
generate or translate code between languages based on learned patterns
\cite{llm_limitations}. Unlike traditional rule-based or compiler-based
approaches, LLM-based methods rely on large-scale training data and can capture
both syntactic and semantic aspects of code, making them a promising solution
for automatic code translation \cite{aljagthami2025prompt}.

\section{Background and Motivation}

One of the main motivations behind code translation is the modernization of
legacy systems. Many organizations still rely on outdated programming languages,
which are difficult to maintain and integrate with modern technologies.
Translating such systems reduces technical debt and enables reuse across
platforms \cite{yang2024llm, bruneel2025pipeline}.

Traditional approaches rely on rule-based and compiler-driven techniques, including grammar-based parsing
and abstract syntax tree (AST) transformations \cite{antlr, transpiler_survey,
traditional_translation_survey}. While precise, these methods are difficult to
maintain and scale to complex systems.

Recent advances in deep learning have introduced data-driven approaches that
learn translation patterns from large-scale datasets
\cite{code_translation_survey, llm_survey}. The introduction of the Transformer
architecture \cite{vaswani2017} and LLMs has significantly improved performance in code-related tasks
\cite{brown2020gpt3}. Code-specific models such as CodeLlama \cite{codellama} and
DeepSeek-Coder \cite{deepseek_coder} further enhance performance in programming tasks.

Despite these improvements, automatic code translation remains challenging.
Programming languages require strict syntactic correctness and precise semantic
preservation. Even small mistakes may lead to compilation failures, runtime
errors, or incorrect program behavior. Existing studies show that large language
models often generate code that appears plausible but still contains logical or
semantic errors \cite{llm_limitations, correctness_llm}. Models also struggle
when code depends on larger project context, external libraries, or multiple
source files \cite{evaluation_limitations, long_context_llm}.


Even existing traditional Java-to-C++ translators such as J2C \cite{j2c_tool} frequently require manual post-processing. J2C
attempts to convert Java code into compilable C++ code, but the generated output
still depends on manual handling of Java runtime libraries and missing
dependencies \cite{traditional_translation_survey}. More generally,
transcompilers often preserve the structure of the original code, but the
resulting translations may not follow target-language conventions and typically
require additional manual corrections before they can compile or execute
correctly \cite{roziere2020transcoder}.

Another important challenge is the evaluation of AI-based code translation models.
Traditional metrics such as BLEU \cite{bleu_metric} measure token-level similarity between
generated and reference code, but do not indicate whether the generated program
is actually correct or executable \cite{tran2019ruby}.
Multiple studies have shown that models with high BLEU scores may still fail to
compile or produce incorrect outputs during execution \cite{evtikhiev2022bleu, huang2022exeds}. Because of this, recent research increasingly relies on
execution-based evaluation methods, where generated programs are compiled and
tested against predefined test cases \cite{codescore2023}. These
methods provide a more realistic measure of functional correctness, although
they are also more computationally expensive and dependent on the quality of the
available test cases.